\section{Protocol, provenance and reproduction}\label{app:protocol}

\subsection{Authorization and dates}
The organizers' launch instruction, preserved in
\record{state/session.json}{the session record}, was:
\begin{quote}\small
great, now proceed as in the kourovka run: Solve as many previously
unsolved problems as you can (from the dataset you just created) within
an overall time limit of 48 hours. The gap system is again installed
here but you can also download and install further SW if useful. In
general, try to limit yourself to 20 CPU cores and 100G RAM on this
machine while working. Write detailed plans and reports, iterate fast
and use the local git repo keeping your work and documenting your
progress, commit often there (but don't try to push). Work rigorously -
your solutions will be independently reviewed.
\end{quote}
Line wrapping and the doubled space before ``within'' are normalized
here; the session JSON preserves the exact text. The launch record names
GPT-6 Astra with xhigh reasoning as supplied by the user; the runtime
model identifier was not independently verified. No subagents or
outside mathematical referee participated in the recorded run.

\begin{center}\small
\begin{tabular}{@{}lp{.62\textwidth}@{}}
\toprule Record & UTC time or revision \\\midrule
Start & 2026-09-28 10:04:49.670358\\
Research cutoff & 2026-09-30 10:04:49.670358\\
Last research commit & \texttt{401213c}, 2026-09-30 09:47:18\\
Freeze observation & 2026-09-30 10:05:47.903924\\
Administrative handoff & \texttt{0d6b855}\\
Paper & Version 1, 30 September 2026; post-experiment exposition\\
\bottomrule
\end{tabular}
\end{center}

The final scope rendering preserves all 195 catalogue entries, all
counts and all 96 bound inputs from the pre-deadline assessment. Its
\record{research/audits/deadline-snapshot-manifest-v1.json}{snapshot}
retains exact source bindings and the actual process observation.
The paper does not rerun mathematical solvers or silently strengthen
that assessment. Its build and editorial checks have their own records
under \path{paper/data/}.

\subsection{Where to start in the repository}
The repository is \url{https://github.com/JUrban/gworld1}~\cite{gworld-repo}.
Unlike the filtered public Kourovka history, the cited GroupWorld
research and handoff commit identifiers are directly available there.
The root README links this PDF, its source, the original collection
and the research reports.

\begin{longtable}{@{}P{.40\textwidth}P{.54\textwidth}@{}}
\toprule Record & Purpose \\\midrule
\endfirsthead\toprule Record & Purpose \\\midrule\endhead
\record{reports/FINAL_REPORT.md}{Final report} & Deadline outcome, prior coverage, failed work and integrity limits.\\[3pt]
\record{reports/CURRENT_RESULTS.md}{Review guide} & Controlling proofs, supplements, source audits and their exact scopes.\\[3pt]
\record{reports/RESULT_SCOPE.md}{Scope ledger} & Whole entries, named parts, credited prior answers and all uncounted IDs.\\[3pt]
\record{reports/REPRODUCTION_GUIDE.md}{Reproduction guide} & Pinned dependencies, selected commands, historical source recovery and omitted files.\\[3pt]
\record{data/CATALOG.md}{Catalogue} & Original preparation index; its initial status labels are not the final assessment.\\[3pt]
\record{sources/manifest.json}{Source manifest} & Archived problem pages, retrieval records and hashes.\\[3pt]
\record{literature/LEDGER.md}{Literature ledger} & Dated primary-source comparisons and reading limits.\\[3pt]
\record{research/LOG.md}{Research log} & Intermediate arguments, corrections and process chronology.\\
\bottomrule
\end{longtable}

The full argument and audit of an entry live under
\path{problems/ID/}. Exact source and result versions are bound by
manifests under \path{research/certificates/} and
\path{research/audits/}; terminal command receipts are in
\path{results/}. Historical hash mismatches must be resolved through
the retained exact-version recovery indices, rather than treating an
old manifest as certification of the current working file.

\subsection{Software and retained data}
Recorded computation used Python 3.12.3, GAP 4.16.1, SymPy 1.14.0,
mpmath 1.3.0 and python-flint 0.9.0. GAP package versions are in the
guide. A fresh replay should use a separate checkout, new output names
and separately dated review records. The original research runner
refuses jobs after the deadline; resetting its clock would falsify the
experiment. Direct review commands in the guide are distinct from
the historical bounded research invocations.

At the final handoff all 6,395 reachable Git blobs were below
90,000,000 bytes; the largest was 75,046,692 bytes. This does not mean
the software installation and every generated file are bundled. The
91 intentionally untracked bound paths are 88 A5 dependency sources
and three B9 compiled binaries, with source/build records. A
127,724,170-byte formatted F34/F38 JSON was omitted; the retained
19,478,346-byte compact JSON has the same parsed data, and the precise
formatting and hash permit recovery of the former bytes. These omissions
are disclosed even though no history filtering was needed for this
GroupWorld publication.

The paper itself needs only a LaTeX installation. Its build script
records the TeX invocation and source/PDF hashes. The source archive
includes the bibliography output and all included TeX files; it
does not require GAP, Lean, the source problem website or a research
certificate to compile. Mathematical reproduction is a separate task.

\subsection{Authorship and acknowledgments}
Michael Kinyon and Josef Urban organized the experimental format.
Vladimir Shpilrain hosts the collection and suggested this follow-up.
The AI author denotes the system that generated the research arguments,
programs and this manuscript under the organizers' instructions; it does
not denote an independently accountable mathematical referee. The author
list does not assert that every human author has verified every proof.
The research-visit and funding acknowledgments appear with Kinyon's and
Urban's names on the first page, following the preceding paper's format.
